\documentclass[11pt]{article}

\usepackage[margin=25mm]{geometry}

\usepackage{amsmath,amssymb,fontspec,unicode-math}

\IfFontExistsTF{DejaVu Serif}{\setmainfont{DejaVu Serif}}{\setmainfont{Latin Modern Roman}}

\IfFontExistsTF{Latin Modern Math}{\setmathfont{Latin Modern Math}}{\setmathfont{DejaVu Math TeX Gyre}}

\usepackage{xurl,hyperref}

\setlength{\parindent}{0pt}\setlength{\parskip}{6pt}

\begin{document}

\title{Online Portfolio Selection with ML Predictions}

\author{ELVIA preliminary AI analysis}\date{21 October 2026}\maketitle

Preliminary assessment: Minor.\par

Independent scientific review remains pending. Author adoption has not been requested.

The fractional-betting lemma puts a zero-growth threshold below one half when the stated formula places it above one half. A ranking notebook also misgroups asset-major rows as date queries. RAM’s geometric-mean floor remains distinct from BCRP universality.

\section*{Review scope}

Latest-version, claim-focused review of the cited author-paper passages and applicable current author-release paths. The stated counterexamples and scope controls were checked analytically; no author program, old verifier, training run or benchmark was executed.

\section*{Problem 1: The fractional-trust threshold lies in the wrong interval}

Source: Lemma A.1(3), PDF pp.20–21; Section 2, PDF p.5

Assessment: supported; impact: minor

\subsection*{Source claim (paraphrase)}

Lemma A.1(3) places the zero-growth forecast error \(\epsilon^*(\lambda)\) in \((0,1/2)\) for its stated fractional-trust range.

\subsection*{Analysis}

At \(\lambda=1/2\), correct and incorrect predictions give wealth factors \(13/8\) and \(7/8\). Chance-level expected log growth is \(\tfrac12\log(91/64)>0\). Since growth decreases linearly in the prediction-error probability, the root is \(\log(13/8)/[\log(13/8)-\log(7/8)]\approx0.78429\), above \(1/2\). More generally the interior range \(1/3<\lambda<1\) has its root in \((1/2,1)\).

\subsection*{Finding and scope}

The below-chance fragility interpretation of this lemma is incorrect. The RAM guarantee is a separate statement.

\subsection*{Proposed correction}

Use the actual logarithmic root and distinguish the endpoints \(\lambda=1/3\) and \(\lambda=1\). Fractional trust can retain positive growth at chance-level forecasts in this two-asset example.

An explicit local correction is described; author adoption and a corresponding published erratum are not established.

\section*{Problem 2: Ranking groups do not match the row order}

Source: Section 3, PDF pp.9–10, Table 4; official notebook cells 34–35

Assessment: supported; impact: minor

\subsection*{Source claim (paraphrase)}

The official notebook’s `prepare\_long` and `build\_lgb\_dataset`, cell 34, are intended to form one cross-sectional ranking query per date.

\subsection*{Analysis}

The melt operation yields asset-major rows, but the supplied group sizes assume contiguous same-date rows. In a complete panel of three assets and four dates, the first size-three query contains the first asset on dates one through three rather than three assets on one date. Labels, features and weights therefore enter the wrong query grouping.

\subsection*{Finding and scope}

The notebook’s ranking interface is inconsistent. The published wealth values are not recalculated or directly contradicted.

\subsection*{Proposed correction}

Sort the entire feature/label/weight table by date before splitting arrays and constructing consecutive query groups. This repairs membership without claiming recovered investment performance.

An explicit local correction is described; author adoption and a corresponding published erratum are not established.

\section*{Problem 3: The geometric-mean floor does not imply BCRP universality}

Source: Theorem 2.1, PDF p.6; introduction/discussion protection interpretation

Assessment: not\_issue; impact: no\_issue\_found

\subsection*{Source claim (paraphrase)}

Theorem 2.1 supplies \(S_n^{\mathrm{RAM}}\geq(\prod_{i,j}x_j(i))^{1/m}\).

\subsection*{Analysis}

With constant two-asset returns \((2,1/2)\) and consistently wrong ranks, the stated RAM behavior cycles to wealth \(1\) at even horizons, whereas the best constant rebalanced portfolio attains \(2^n\). This separates the geometric-mean comparator from a sublinear log-regret guarantee against BCRP.

\subsection*{Finding and scope}

Theorem 2.1 remains valid in this example. It should be interpreted with its proved comparator.

\section*{Limits}

Preliminary AI assessment; independent human scientific adjudication remains pending.\par

An analytic counterexample or a release discrepancy does not by itself remeasure the paper’s reported experiments.\par

\section*{Source identity}

NeurIPS 2025 official proceedings PDF (current publisher bytes, checked 2026-10-21)

\url{https://proceedings.neurips.cc/paper_files/paper/2025/file/76981cb7a06b5ad53639f5bbe9924388-Paper-Conference.pdf}

Reviewed PDF SHA-256: \nolinkurl{909d36f26e2aa1f04c0059b9733c93ba3c2d7c9ca6421183d930f35c00083ade}

The current official publisher PDF is the same work as the original case: title, author identity and publisher hash URL agree, and the freshly downloaded scientific PDF is byte-identical to the prior source. The actual download timestamp and SHA-256 are recorded for Issue \#27. Exact-title arXiv search did not return a matching later version; background references were excluded.

\end{document}
